\chapter{Model Architectures}
\label{ch:models}

This chapter describes the KAN layer, the two recurrent baselines, and the
places where KAN layers are inserted into them. Every equation corresponds
directly to the code in \texttt{kanrec/kan.py}, \texttt{kanrec/session.py} and
\texttt{kanrec/basket.py}.

\section{The KAN Layer}

\subsection{Definition}

A KAN layer with $d_{in}$ inputs and $d_{out}$ outputs computes
\begin{equation}
y_o = \sum_{i=1}^{d_{in}} \Big( w^{b}_{o,i}\, \mathrm{silu}(x_i) + \sum_{k=1}^{K} c_{o,i,k}\, B_k(x_i) \Big) + b_o,
\qquad o = 1, \ldots, d_{out},
\label{eq:kan}
\end{equation}
where $B_1, \ldots, B_K$ are fixed basis functions, $c_{o,i,k}$ are learnable
coefficients and $w^b_{o,i}$, $b_o$ form a linear ``base'' path on
$\mathrm{silu}(x)$, as in the efficient implementations of the original KAN
\cite{liu2025kan}. Each of the $d_{in} d_{out}$ edges therefore carries its own
learnable univariate function
\begin{equation}
\phi_{o,i}(x) = w^b_{o,i}\,\mathrm{silu}(x) + \sum_{k} c_{o,i,k} B_k(x),
\end{equation}
which can be plotted for interpretation (Section~\ref{sec:interp}). The layer
has $d_{in} d_{out} (K + 1) + d_{out}$ parameters, compared with
$d_{in} d_{out} + d_{out}$ for a linear layer.

\subsection{Basis Families}

Four basis families are implemented:

\begin{itemize}
    \item \textbf{B-spline} (cubic, $p = 3$, grid size $G$). With uniform knots
          $t_j = t_0 + jh$ on an extended grid that adds $p$ knots beyond each end
          of the range, the bases follow the Cox-de~Boor recursion
          \begin{align}
          B^{(0)}_j(x) &= \mathbb{1}[t_j \le x < t_{j+1}], \\
          B^{(q)}_j(x) &= \frac{x - t_j}{t_{j+q} - t_j}\, B^{(q-1)}_j(x)
                       + \frac{t_{j+q+1} - x}{t_{j+q+1} - t_{j+1}}\, B^{(q-1)}_{j+1}(x),
          \end{align}
          giving $K = G + p$ bases that sum to one everywhere inside the grid.
          The unit tests check this partition of unity and compare the bases with
          the closed-form uniform cubic B-spline.
    \item \textbf{Gaussian RBF} (FastKAN \cite{li2024fastkan}):
          $B_k(x) = \exp\!\big(-((x - \mu_k)/h)^2\big)$ with $K = G$ centres
          $\mu_k$ evenly spaced over the grid range and $h$ their spacing.
    \item \textbf{Chebyshev} \cite{ss2024cheby}: $B_k(x) = T_{k-1}(\tanh x)$,
          $k = 1, \ldots, G+1$, where $T_n$ are Chebyshev polynomials,
          computed with $T_0 = 1$, $T_1(u) = u$, $T_{n}(u) = 2uT_{n-1}(u) - T_{n-2}(u)$.
    \item \textbf{Group rational} (GR-KAN \cite{yang2025kat}): the inputs are split
          into 8 groups; within a group every input passes through the same
          learnable rational function
          \begin{equation}
          R(x) = \frac{\sum_{m=0}^{5} a_m x^m}{1 + \big|\sum_{n=1}^{4} b_n x^n\big|},
          \end{equation}
          followed by a linear layer. Functions are shared per group rather than
          per edge, so this variant is much smaller.
\end{itemize}

The B-spline and RBF grids span $[-2, 2]$, which covers inputs normalised by
LayerNorm; the default grid size is $G = 8$.

\subsection{Residual Block and MLP Control}

KAN layers are inserted as residual blocks,
\begin{equation}
\mathrm{Block}(\mathbf{h}) = \mathbf{h} + f\big(\mathrm{LayerNorm}(\mathbf{h})\big),
\label{eq:block}
\end{equation}
with $f$ a $d \to d$ KAN layer. The output layer of $f$ is initialised to zero,
so every block starts as the identity and every variant starts as exactly its
baseline model; the LayerNorm keeps the inputs inside the basis grid. For each KAN block, a \textbf{control block} uses the same
wrapper with $f(\mathbf{x}) = W_2\, \mathrm{silu}(W_1 \mathbf{x} + b_1) + b_2$, a
two-layer MLP whose hidden width is chosen so that it has the same number of
parameters as the KAN (within 2\%). A KAN block that beats the baseline but not
its control has shown only that the extra capacity helps.

\section{Session Model: GRU4Rec}

GRU4Rec \cite{hidasi2016,hidasi2018} with the official settings for YooChoose
uses one GRU layer of 480 units and \emph{constrained embeddings}: the item
embedding matrix $E$ is shared between input and output. For an input item $i_t$,
\begin{align}
\mathbf{x}_t &= E_{i_t}, \\
\mathbf{r}_t &= \sigma(W_{ir}\mathbf{x}_t + b_{ir} + W_{hr}\mathbf{h}_{t-1} + b_{hr}), \\
\mathbf{z}_t &= \sigma(W_{iz}\mathbf{x}_t + b_{iz} + W_{hz}\mathbf{h}_{t-1} + b_{hz}), \\
\mathbf{n}_t &= \tanh\big(W_{in}\mathbf{x}_t + b_{in} + \mathbf{r}_t \odot (W_{hn}\mathbf{h}_{t-1} + b_{hn})\big), \\
\mathbf{h}_t &= (1 - \mathbf{z}_t) \odot \mathbf{n}_t + \mathbf{z}_t \odot \mathbf{h}_{t-1},
\end{align}
and the score of item $j$ is $s_j = E_j^\top \mathbf{h}_t + \beta_j$.
(These are the PyTorch \texttt{GRUCell} equations; the first draft's candidate
equation was incorrect.) Training uses session-parallel mini-batches of 48
sessions, 2048 additional negative items shared across the mini-batch and
sampled in proportion to popularity$^{0.2}$, a softmax cross-entropy loss over
the targets and samples with a $\log Q$ popularity correction, dropout 0.2 on the
hidden state, and Adagrad with learning rate 0.07.

\section{KAN Placements in GRU4Rec}

\begin{figure}[H]
\centering
\resizebox{\textwidth}{!}{%
\begin{tikzpicture}[
  node distance=0.55cm and 0.9cm,
  box/.style={draw, rounded corners, minimum height=0.7cm, minimum width=1.9cm, align=center, font=\small},
  kan/.style={box, fill=orange!20},
  arr/.style={-{stealth}, thick}]
\node[box] (item) {item $i_t$};
\node[box, right=of item] (emb) {embedding\\$E_{i_t}$};
\node[kan, right=of emb] (kin) {KAN input\\block (S2)};
\node[box, right=of kin, fill=blue!8] (gru) {GRU cell\\(S3: KAN-GRU)};
\node[kan, right=of gru] (kout) {KAN head\\block (S1)};
\node[box, right=of kout] (score) {scores\\$E\mathbf{h}' + \beta$};
\draw[arr] (item) -- (emb);
\draw[arr] (emb) -- (kin);
\draw[arr] (kin) -- (gru);
\draw[arr] (gru) -- (kout);
\draw[arr] (kout) -- (score);
\draw[arr] ([xshift=0.4cm]gru.south) .. controls +(0.6,-1.0) and +(-0.6,-1.0) .. node[below, font=\scriptsize]{$\mathbf{h}_{t-1}$} ([xshift=-0.4cm]gru.south);
\end{tikzpicture}}
\caption{KAN placements in GRU4Rec. Orange blocks are optional residual KAN (or
MLP control) blocks; the GRU cell can be replaced by the KAN-GRU cell.}
\label{fig:placements}
\end{figure}

\begin{itemize}
    \item \textbf{S1, KAN head}: $\mathbf{h}'_t = \mathrm{Block}(\mathbf{h}_t)$
          before scoring. The block refines the session representation; scoring
          stays a dot product, so sampled negatives and constrained embeddings
          work unchanged. (A KAN on the item-scoring layer itself would need
          $480 \times 17{,}000 \times K$ coefficients.)
    \item \textbf{S2, KAN input}: $\mathbf{x}'_t = \mathrm{Block}(E_{i_t})$, a
          non-linear transformation of the item embedding before the GRU.
    \item \textbf{S1+S2}: both blocks.
    \item \textbf{S3, KAN-GRU cell}: the update and reset gates keep their linear
          + sigmoid form, as in TKAN \cite{genet2024}, and only the candidate
          state is computed by a KAN:
          \begin{equation}
          \mathbf{n}_t = \tanh\big(\mathrm{KAN}([\mathbf{x}_t;\ \mathbf{r}_t \odot \mathbf{h}_{t-1}])\big).
          \end{equation}
    \item \textbf{Draft models}, for continuity with the first draft:
          \emph{Draft Pure KAN} replaces the GRU with the ungated cell
          $\mathbf{h}_t = \tanh(\mathrm{LayerNorm}(\mathrm{KAN}([\mathbf{x}_t;\mathbf{h}_{t-1}])))$
          with a correct cubic B-spline basis, and \emph{Draft Hybrid} stacks that
          cell on top of a GRU layer.
\end{itemize}

S1 and S2 are paired with parameter-matched MLP controls (\emph{+MLP head},
\emph{+MLP input}), and S3 with the \emph{MLP-GRU cell}, whose candidate state
$\tanh(\mathrm{MLP}([\mathbf{x}_t; \mathbf{r}_t \odot \mathbf{h}_{t-1}]))$ uses an MLP with
the same number of parameters as the KAN. Draft Hybrid has two recurrent layers,
so it is compared with a two-layer GRU4Rec (same depth) and with \emph{Draft
Hybrid (MLP)}, in which the KAN of the ungated cell is replaced by a
parameter-matched MLP. Unless stated otherwise the KAN layers use the RBF basis
with $G=8$; the other bases are compared in an ablation.

\section{Next-Basket Model: Basket-GRU}

The basket-GRU of van Maasakkers et al.\ \cite{vanmaasakkers2023} reads the
multi-hot vector $\mathbf{m}_t \in \{0,1\}^N$ of basket $t$ and outputs a purchase
probability for every one of the $N$ items:
\begin{align}
\mathbf{x}_t &= W_{in}\mathbf{m}_t, \\
\mathbf{h}_t &= \mathrm{GRU}(\mathbf{x}_t, \mathbf{h}_{t-1}), \\
\hat{\mathbf{y}}_{t+1} &= \sigma(W_{out}\mathbf{h}_t + \mathbf{b}),
\end{align}
trained with binary cross-entropy against $\mathbf{m}_{t+1}$, averaged over all
items and over every step of every training history. The output bias $b_j$ is
initialised to the log-odds of item $j$'s share of training baskets, so that
training starts from the popularity predictor; without this, the biases need
thousands of updates to move from probability 0.5 towards the true base rates
of roughly $10^{-3}$, and short training runs stay at chance level. $W_{in}\mathbf{m}_t$ is
implemented as a sum-pooled embedding bag. Sequences of different lengths are
padded on the right and the representation of a user is read at their true last
basket (not at the last padded position, the error suspected in the first draft).
The KAN placements mirror those of GRU4Rec: \emph{+KAN head} applies
Equation~\eqref{eq:block} to $\mathbf{h}_t$, \emph{+KAN input} to $\mathbf{x}_t$,
and \emph{KAN-GRU cell} uses the cell of S3, each with its MLP control. As before, no KAN is applied on the
$N$-dimensional side.

\section{KAN-NBR: An Interpretable Repeat-Aware Model}
\label{sec:kannbr}

Because repeat purchases account for most of the achievable next-basket
accuracy \cite{li2023}, we also test a KAN where its interpretability can matter:
on a few meaningful scalar features rather than on hundreds of hidden units.
For user $u$ and candidate item $j$, where the candidates are the items $u$ has
bought before plus the 100 most popular items, KAN-NBR computes seven features
from the history $B^u_1, \ldots, B^u_n$:

\begin{enumerate}
    \item $\log(1 + \text{baskets since the last purchase of } j)$;
    \item $\log(1 + c_{uj})$, where $c_{uj}$ is the number of baskets containing $j$;
    \item the purchase share $c_{uj}/n$;
    \item the decayed frequency $\sum_t 0.9^{\,n-t}\, \mathbb{1}[j \in B^u_t]$;
    \item the log popularity of $j$ across all users;
    \item $\log(1 + \text{overdue ratio})$, the baskets since the last purchase
          divided by the user's mean interval between purchases of $j$;
    \item whether $j$ is in the history at all.
\end{enumerate}

After batch normalisation, a single KAN layer maps the seven features to a
purchase logit:
\begin{equation}
\mathrm{logit}\, P(j \in B^u_{n+1}) = \sum_{f=1}^{7} \phi_f(z_f) + b,
\end{equation}
which is a generalised additive model whose shape functions $\phi_f$ are learned
KAN edge functions and can be plotted directly. A two-layer KAN
$[7 \to 4 \to 1]$ allows interactions between features. Controls are logistic
regression on the same features (linear $\phi_f$) and MLPs with the same number
of parameters as each KAN. Training pairs are built from the last three positions
of every user's history (predicting basket $t$ from the baskets before it), so
the held-out final basket is never seen in training. The model is trained with
binary cross-entropy and Adam, with early stopping on validation Recall@10.

\section{Baselines}

\textbf{Session}: \emph{Pop} (global popularity), \emph{S-Pop} (items already in
the session by in-session count, then popularity) and \emph{Item-kNN}
\cite{hidasi2016}, which scores item $j$ by its similarity to the current item $i$,
$\mathrm{sim}(i,j) = n_{ij} / (n_i^{0.5} n_j^{0.5} + 20)$, where $n_{ij}$ counts
sessions containing both items.

\textbf{Next-basket}: \emph{G-TopFreq} (global popularity), \emph{P-TopFreq}
(the user's own purchase frequencies), \emph{GP-TopFreq} (personal frequencies,
remaining positions filled by global popularity), \emph{Last basket}, and
\emph{TIFU-KNN} \cite{hu2020}. TIFU-KNN represents a user by time-decayed item
frequencies: the history is cut into groups of $m$ baskets, baskets within a
group are weighted by $r_b^{\,\text{age in group}}$ and groups by
$r_g^{\,\text{group age}}$. The prediction is
$\alpha \mathbf{u} + (1-\alpha)\,\overline{\mathbf{u}}_{\mathrm{kNN}}$, mixing the user's
vector with the mean vector of their $k$ most similar users. Its hyperparameters
are tuned on validation users.
